\chapter{The Serial Test}

\section{Intuition from Dice}
The frequency test asks whether each outcome appears often enough in the long
run.  But a loaded process can hide in a different way: each face might appear
with the \emph{right overall frequency}, yet the results might not be
\emph{independent}.

Here is a simple example that makes the point without any statistics.  In the
real world, ``sunrise'' and ``sunset'' occur with the same overall frequency:
one of each per day.  But that does not make the sequence random.  In fact it
is almost perfectly predictable: a sunrise is followed by a sunset, and a
sunset is followed by a sunrise.  The \emph{frequencies} are balanced, yet the
\emph{order} is highly structured.

The \textbf{serial test} targets exactly this kind of problem.  It asks:
\begin{quote}
  \emph{Do consecutive outcomes behave like independent draws, or do certain
  short patterns (pairs, triples, \dots) occur too often or too rarely?}
\end{quote}

\section{What the Serial Test Checks}
Instead of counting single outcomes, the serial test counts \emph{tuples}.
For example:
\begin{itemize}
  \item a \textbf{2-tuple} (pair) test counts how often \((a,b)\) occurs in two
        consecutive positions;
  \item a \textbf{3-tuple} (triple) test counts how often \((a,b,c)\) occurs in
        three consecutive positions.
\end{itemize}

If the mechanism is fair \emph{and independent}, then the probability of a tuple
is the product of the individual probabilities.  For a fair dice, for example,
\[
  P(1,1) = P(1)\cdot P(1) = \frac{1}{6}\cdot\frac{1}{6} = \frac{1}{36}.
\]
The serial test compares the observed tuple counts to these product
probabilities.

\section{A PowerSpin Example: the 3-Wheel Serial Test}
PowerSpin has \(27\) positions on the wheel: numbers \(1\)--\(24\) and 3 special
symbols.  In the code below we treat the 3 symbols as one merged category
called \texttt{symbol}.  This gives \(25\) categories in total:
\[
  \{1,2,\dots,24,\texttt{symbol}\}.
\]

The model used by the code is:
\begin{itemize}
  \item per wheel, each number \(1\)--\(24\) has probability \(1/27\);
  \item per wheel, \texttt{symbol} has total probability \(3/27\).
\end{itemize}

If the three wheels are independent, the probability of a triple
\((a,b,c)\) is the product of the three wheel probabilities.  For example:
\[
  P(\texttt{symbol},\texttt{symbol},\texttt{symbol}) = \left(\frac{3}{27}\right)^3 = \frac{1}{729}.
\]

The serial test then counts how often each of the \(25\times 25\times 25\)
triples occurs in the data and runs a chi-squared goodness-of-fit test against
this product distribution (Section~\ref{sec:chi-squared-test}).

\section{A Practical Warning: Too Many Categories}
This 3-wheel serial test is conceptually clean, but it creates a very large
number of categories:
\[
  K = 25^3 = 15\,625.
\]

With only one month of data, the \emph{expected} count per category is often
small.  For example, with \(N = 22\,320\) triples, the average expected count is
about \(22\,320/15\,625 \approx 1.43\).  This matters because the usual
chi-squared approximation is most reliable when expected counts are not too
small (Section~\ref{sec:chi-squared-test}).

In practice, when expected counts are tiny, you have three options:
\begin{itemize}
  \item collect much more data;
  \item reduce the number of categories (for example, test 2-tuples instead of
        3-tuples, or merge outcomes into coarser bins);
  \item use a simulation-based calibration (generate many synthetic datasets
        from the model and compare your statistic to that empirical reference).
\end{itemize}

\section{How to Perform the Test}
\begin{enumerate}
  \item \textbf{Normalize categories.}\\
        Parse the CSV values.  If a value is an integer from 1 to 24, treat it
        as that number; otherwise treat it as \texttt{symbol}.  (This matches
        \texttt{normalize\_category} in the code.)

  \item \textbf{Convert categories to indices.}\\
        Map \(1\)--\(24\) to indices \(0\)--\(23\), and \texttt{symbol} to index
        \(24\).  (This matches \texttt{category\_index}.)

  \item \textbf{Count observed triples.}\\
        For each row that contains all three wheel results, increment a
        3-dimensional array \(\texttt{observed}[i,j,k]\).  At the end you have:
        \begin{itemize}
          \item a count table of shape \(25\times 25\times 25\);
          \item the number of valid triples \(N\) (rows with all 3 results).
        \end{itemize}

  \item \textbf{Build the expected triple probabilities.}\\
        Create a length-25 probability vector
        \(\texttt{p} = (1/27,\dots,1/27,3/27)\).  Then compute the product
        distribution
        \[
          \texttt{expected}[i,j,k] = \texttt{p}[i]\cdot \texttt{p}[j]\cdot \texttt{p}[k].
        \]
        (This matches \texttt{expected\_triple\_distribution}.)

  \item \textbf{Run the chi-squared test and get a $p$-value.}\\
        Feed \texttt{observed} and \texttt{expected} into the chi-squared
        procedure to obtain a single $p$-value.  A typical implementation uses
        \(\nu = K-1\) degrees of freedom, where \(K\) is the number of outcome
        categories (here \(K = 25^3\)), and computes the upper-tail probability
        from the chi-squared distribution.
\end{enumerate}

\section{One-Year Result}
Using one year of PowerSpin data, the 3-wheel serial chi-squared test produced a
$p$-value of \(0.911716\).  This is \emph{somehow suspicious}. Further tests from a broader battery will help determine whether this
is simply good luck in the sample or whether the mechanism is not behaving like
genuine randomness.




